\section{Mathematical background}


\subsection{The forward operator}
\input{figures/forward_operator.tex}
Let $\Omega \subset \mathbb R^2$ be a bounded domain with boundary $\partial \Omega$, and consider the boundary value problem (BVP) in Equation~\eqref{eq:pde}. The coefficient $\sigma = \sigma(x) > 0$ denotes the diffusion coefficient and determines the rate and direction of diffusion in the problem. It may be either scalar-valued or matrix-valued. In the latter case, we assume that $\sigma(x)$ is symmetric and positive definite (see Appendix~\ref{sec:positive_definiteness}) to keep the problem elliptic. The reaction coefficient $c = c(x) > 0$ describes the absorption or decay within the problem. Finally, $f = f(x)$ denotes the source term, which is the unknown input of the inverse problem.
\nomenclature[A]{$\Omega$}{Computational domain}
\nomenclature[A]{$\partial\Omega$}{Boundary of the domain}
\nomenclature[B]{$\sigma,\, \sigma(x)$}{Diffusion coefficient}
\nomenclature[B]{$c,\, c(x)$}{Reaction coefficient}

Equation~\eqref{eq:pde} is a linear diffusion-reaction problem with Neumann boundary conditions, and defines a mapping
\begin{equation}\label{eq:S_cont}
    \mathcal S: L^2(\Omega) \rightarrow H^1(\Omega), \quad \mathcal{S} f = u.
\end{equation}
Here, $L^2(\Omega)$ is the space of square integrable functions and $H^1(\Omega) \subset L^2(\Omega)$ the first-order Sobolev space, see Appendix~\ref{appendix:function_spaces}
\nomenclature[C]{$\mathcal{S}$}{PDE operator}
\nomenclature[D]{$L^2(\Omega)$}{Space of square integrable functions}
\nomenclature[D]{$H^1(\Omega)$}{First-order Sobolev space}

To model boundary observations, we introduce the trace operator
\begin{equation}\label{eq:trace}
    \mathcal T: H^1(\Omega) \rightarrow L^2(\partial\Omega), \quad \mathcal{T} u = u|_{\partial\Omega}
\end{equation}
\nomenclature[C]{$\mathcal{T}$}{Trace operator}
The term $\mathcal T u$ represents the observed values of $u$ on the boundary of $\Omega$.

The forward operator is then defined as the composition
\begin{equation}\label{eq:K}
    \mathcal K = \mathcal T \circ \mathcal S, \quad \mathcal{K} f = u|_{\partial\Omega},
\end{equation}
\nomenclature[C]{$\mathcal{K}$}{Forward operator}
which is a mapping from the source function $f$ to the observed output $u|_{\partial\Omega} = \mathcal K f$ on the boundary $\partial \Omega$. See Figure~\ref{fig:forward_operator}.

The associated inverse problem can now be formulated as follows: given boundary data $u|_{\partial\Omega}$, determine a source function $f$ such that
\begin{equation}
    \mathcal{K} f = u|_{\partial\Omega}.
\end{equation}
Solving this problem computationally requires discretizing $\mathcal{K}$, understanding the ill-posedness of the resulting system, and regularizing it to obtain stable solutions, which we address in the following sections.


\subsubsection{Ill-posedness}
A problem is said to be well-posed if its solution exists, is unique, and is stable in the sense that it depends continuously on the data, as formulated by Hadamard \cite{hansen2010discrete}. The inverse problem $\mathcal{K}f = y$ is ill-posed.

The null space of $\mathcal{K}$ is non-trivial: any $f_0$ satisfying $\mathcal{K}f_0 = 0$ can be added to a solution without changing the boundary data $y$. The solution is therefore not unique.

The instability of the solution $f$ is more severe in practice. The compactness of $\mathcal{K}$, and the ellipticity of the PDE in Equation~\eqref{eq:pde}, results in a rapid decay of its singular values towards zero \cite{evansPartialDifferentialEquations2010,engl1996regularization}. Solution components corresponding to small singular values are strongly suppressed by the forward operator and cannot be recovered from noisy data: a small perturbation in $y$ produces a large error in the reconstructed source. This motivates the regularized formulation stated in Equation~\eqref{eq:problem}.


\subsection{Finite element discretization}
\subsubsection{The finite element method}
\input{figures/triangulation.tex}
Consider the PDE in Eq.~\eqref{eq:pde}. Multiplying the equation by a test function $v \in H^1(\Omega)$, integrating over $\Omega$, and applying Green's identity yields the weak formulation: find $u \in H^1(\Omega)$ such that
\begin{equation}\label{eq:variational_formulation}
    \int_{\Omega} (\sigma \nabla u \cdot \nabla v + c u v) \, dx = \int_{\Omega} f v \space \, dx,
    \text{ for all } v \in H^1(\Omega).
\end{equation}
For details on the derivation, see Section~4 in Larson and Bengzon \cite{larsonFiniteElementMethod2013}.

Let $\mathcal M$ be a triangulation of $\Omega$ made up of $N$ nodes, with $N_b$ boundary nodes, see Fig.~\ref{fig:triangulation} as an example. Introduce the finite element space $V_h \subset H^1(\Omega)$,
\begin{equation}\label{eq:fem_space}
    V_h = \{v \in C^0(\Omega) : v|_M \text{ is linear for all } M \in \mathcal M \}.
\end{equation}
Let $\{\phi_1, \dots, \phi_N\}$ denote a basis for $V_h$, so that
\begin{equation}
    V_h = \operatorname{span}\{ \phi_1, \dots, \phi_N \}.
\end{equation}
The basis functions $\phi_i$ are piecewise linear hat functions associated with the nodes of the mesh (see \cite[p.~51]{larsonFiniteElementMethod2013}).

\nomenclature[A]{$\mathcal{M}$}{Triangulation (mesh)}
\nomenclature[A]{$N$}{Number of mesh nodes}
\nomenclature[A]{$N_b$}{Number of boundary mesh nodes}
\nomenclature[D]{$V_h$}{Finite element space}

The finite element approximation of \eqref{eq:variational_formulation} then reads: 
find $u_h \in V_h$ such that
 \begin{equation}\label{eq:fem_formulation}
    \int_{\Omega} (\sigma \nabla u_h \cdot \nabla v + c u_h v) \, dx = \int_{\Omega} f v \space \, dx,
    \text{ for all } v \in V_h.
\end{equation}


\subsubsection{The discrete forward operator}
To obtain a discrete representation of the forward operator 
$\mathcal{K} = \mathcal{T} \circ \mathcal{S}$, 
we discretize both the PDE operator $\mathcal{S}$ and the trace operator $\mathcal{T}$ using the finite element space $V_h$.

Let $f_h, u_h \in V_h$ denote the finite element approximations of $f$ and $u|_{\partial\Omega}$, respectively. Since $V_h = \operatorname{span}\{\phi_1,\dots,\phi_N\}$, they can be written as
\begin{equation}\label{eq:uh_and_fh}
    f_h = \sum_{j=1}^{N} x_j \phi_j, \quad
    u_h = \sum_{j=1}^{N} u_j \phi_j, 
\end{equation}
with coefficient vectors $\v u = (u_1,\dots,u_N)^T$ 
and $\v x = (x_1,\dots,x_N)^T$.

Inserting \eqref{eq:uh_and_fh} into the finite element formulation \eqref{eq:fem_formulation} and testing with the basis functions 
$\{\phi_i\}_{i=1}^N$ yields the linear system
\begin{equation}\label{eq:matrix_system}
    A \mathbf{u} = M \mathbf{x},
\end{equation}
where $A \in \mathbb{R}^{N \times N}$ is the stiffness matrix and $M \in \mathbb{R}^{N \times N}$ the mass matrix, see Appendix~\ref{appendix:stiffness}.
\nomenclature[E]{$M$}{Volume mass matrix}%
Equation~\eqref{eq:matrix_system} implies the discrete PDE operator $S$ takes the form
\begin{equation}
    S = A^{-1} M \in \mathbb{R}^{N \times N}
\end{equation}
Note that $S$ is a mapping of coefficients vectors: mapping $\v x$ to $\v u$.

The trace operator is discretized by a matrix 
$T \in \mathbb{R}^{N_b \times N}$ that extracts the boundary coefficients. If $\mathbf{y}$ denotes the vector of boundary coefficients, then
\begin{equation}
    \v y = T \v u.
\end{equation}

Combining the discrete PDE and trace operators yields the matrix representation of the forward operator,
\begin{equation}\label{eq:discrete_K}
    K = T S = T A^{-1} M \in \mathbb{R}^{N_b \times N}.
\end{equation}
\nomenclature[E]{$K$}{Discrete forward operator}%


\subsubsection{Discrete inner product and norm}
Let $V_h$ be the finite element space introduced in Eq.~\eqref{eq:fem_space} and let $f_h, g_h \in V_h$, with coefficient vectors $\v x, \v y \in \mathbb{R}^N$, respectively. The $L^2(\Omega)$ inner product (see Eq.~\eqref{eq:L2_inner_product}) of $f_h$ and $g_h$, can be written as
\begin{equation}
    \langle f_h, g_h \rangle_{L^2(\Omega)} = \v x^T M \v y,
\end{equation}
where $M$ is the mass matrix. 

This induces a discrete inner product on $\mathbb{R}^N$, with associated norm, defined by
\begin{equation}\label{eq:fem_norm}
    \langle \v x, \v y \rangle_M = \v x^T M \v y,
    \quad
    \| \v x \|_M = \sqrt{\v x^T M \v x}.
\end{equation}

Similarly, let $\v u, \v v \in \mathbb{R}^{N_b}$ 
denote coefficient vectors corresponding to the traces 
$u_h|_{\partial\Omega}$ and $v_h|_{\partial\Omega}$. The $L^2(\partial\Omega)$ inner product induces a discrete inner product and norm on $\mathbb{R}^{N_b}$, given by
\begin{equation}\label{eq:fem_norm_ds}
    \langle \mathbf{u}, \mathbf{v} \rangle_{M_{\partial}}
    = \mathbf{u}^T M_{\partial} \mathbf{v},
    \quad
    \| \v  u \|_{M_{\partial}}
    =
    \sqrt{ \v  u^T M_{\partial} \v u},
\end{equation}
where $M_\partial$ denotes the boundary mass matrix, see Appendix~\ref{appendix:stiffness}.
\nomenclature[E]{$M_\partial$}{Boundary mass matrix}%


\subsubsection{The adjoint operator}
Let $\mathcal{K}: L^2(\Omega) \to L^2(\partial\Omega)$ be the bounded linear operator defined in Equation~\eqref{eq:K}. The adjoint operator $\mathcal{K}^*: L^2(\partial\Omega) \to L^2(\Omega)$ is the bounded linear operator satisfying
\begin{equation}\label{eq:K_star_cont}
    \langle \mathcal K f,\, g \rangle_{L^2(\partial\Omega)} =
    \langle f,\, \mathcal K^* g \rangle_{L^2(\Omega)},
\end{equation}
for all $f \in L^2(\Omega)$ and all $g \in L^2(\partial\Omega)$. Such an operator always exists for a bounded linear operator $\mathcal K$ mapping between two Hilbert spaces, see Section~3.9 in Kreyszig for a proof \cite{kreyszigIntroductoryFunctionalAnalysis1989}.

Using the definition in Equation~\eqref{eq:K_star_cont}, we get the weak formulation of $K^*$ (the adjoint problem): find $w \in H^1(\Omega)$ such that
\begin{equation}
    \int_\Omega (\sigma^T \nabla w \cdot \nabla v + c w v) \, dx =
    \int_{\partial\Omega} g v \, ds,
    \text{ for all } v \in H^1(\Omega).
\end{equation}
The solution $w$ then satisfies $\mathcal K^* g = w$. For details, see \textcolor{red}{(???)}.

In the discrete case, under the inner products and norms defined in Equation~\eqref{eq:fem_norm} and \eqref{eq:fem_norm_ds}, the adjoint operator $K^* \in \mathbb{R}^{N \times N_b}$ of $K \in \mathbb{R}^{N_b \times N}$ satisfies
\begin{equation}
    (K \v x)^T M_\partial \v y = \v x^T M K^* \v y,
\end{equation}
for all $\v x \in \mathbb R^N$ and all $\v y \in \mathbb R^{N_b}$. 

Inserting the explicit representation of $K$ given in Equation~\eqref{eq:discrete_K} yields the matrix representation of $K^*$,
\begin{equation}
    K^* = A^{-T} T^T M_\partial,
\end{equation}
where $A$ is the stiffness matrix, $T$ the discrete trace operator, and $M_\partial$ the boundary mass matrix.

\subsection{Tikhonov regularization}

Let $A \in \mathbb{R}^{m \times n}$ and consider the inverse problem $Ax = y$. When the problem is ill-posed or $A$ is ill-conditioned, small perturbations in $y$ cause large changes in the solution. Tikhonov regularization stabilizes the problem by solving
\begin{equation}\label{eq:tikhonov}
    \min_{x \in \mathbb{R}^n}
    \|Ax - y\|^2 + \lambda^2 \|Wx\|^2,
\end{equation}
where $\lambda \ge 0$ is the regularization parameter and $W \in \mathbb{R}^{n \times n}$ is the weight matrix. $W$ can include prior knowledge about the solution to punish specific types of unwanted solutions \cite{hansen2010discrete}. With $W = I$, this is standard Tikhonov regularization.


\subsection{The discrete inverse problem}
Replacing $\mathcal{K}$ by $K \in \mathbb{R}^{N_b \times N}$ and the $L^2$ norms by the FEM inner products in Equations~\eqref{eq:fem_norm} and~\eqref{eq:fem_norm_ds}, the Tikhonov problem takes the discrete form
\begin{equation}\label{eq:discrete_inverse_problem}
    \min_{\mathbf{x} \in \mathbb{R}^N}
    \| K\mathbf{x} - \mathbf{y} \|_{M_\partial}^2
    + \lambda^2 \| W\mathbf{x} \|_M^2,
\end{equation}
where $\mathbf{x} \in \mathbb{R}^N$ is the coefficient vector of $f_h \in V_h$, $\mathbf{y} \in \mathbb{R}^{N_b}$ is the observed boundary data, and $\lambda \geq 0$ is the regularization parameter.

Standard Tikhonov ($W = I$) does not work here. Because $K$ has a large nullspace, the solutions it produces concentrate near the domain boundary regardless of where the true source is located \cite{elvetunRegularizationOperatorSource2020}. Elvetun and Nielsen address this by choosing weights that penalize boundary-concentrated solutions,
\begin{equation*}
    w_i = \|K^+ K e_i\|,
\end{equation*}
where $e_i$ is the $i$th standard basis vector of $\mathbb{R}^N$ \cite{elvetunRegularizationOperatorSource2020}, and $W = \operatorname{diag}(w_1, \dots, w_N)$. In $V_h$ with the inner product induced by $M_\partial$, the weights take the form
\begin{equation}\label{eq:weights_fem}
    w_i =
    \frac{\sqrt{\mathbf{y}_i^T M_\partial \mathbf{y}_i}}{\int_\Omega \phi_i \, dx},
    \quad
    \mathbf{y}_i = K^+ K e_i.
\end{equation}


\subsection{Randomized SVD}
The rSVD is a technique used to compute a rank-$k$ approximation of an $m\times n$ matrix without computing the SVD. The core idea is to use random sampling to find a low-rank matrix $Q$ which approximates the range of $A$, project $A$ onto $\operatorname{span}(Q)$, and compute the SVD of the projection. The projection $B=Q^TA$ will, with high probability, capture the dominant actions of $A$, and the SVD of $B$ will yield a good approximation \cite{halkoFindingStructureRandomness2011}. This procedure is summarized in Algorithm \ref{alg:rsvd}. The construction of an orthonormal basis $Q$ of $Y$ is done through a QR factorization. Note that the approximation $A\approx QQ^T A$ and the approximate factorization $A\approx U\Sigma V^T$ produced by rSVD have the same error. This can be seen by observing that $QQ^TA = QB = Q\tilde U\Sigma V^T$, and setting $U=Q\tilde U$.
\begin{algorithm}[!htbp]
\caption{rSVD}
\label{alg:rsvd}
\begin{algorithmic}[1]
    \Require $A\in\mathbb R^{m\times n}$, target rank $k$, oversampling parameter $p$
    \Ensure Approximate rank-$k$ factorization $U\Sigma V^T$
    \State Generate a random $n\times(k+p)$ test matrix $\Psi$
    \State Compute $Y = A\Psi$
    \State Form an orthonormal basis $Q$ of $Y$
    \State Set $B = Q^TA$
    \State Compute the SVD $B = \tilde U\Sigma V^T$
    \State Set $U = Q\tilde U$
    \State Truncate to rank $k$: $U, \Sigma, V^T \leftarrow \operatorname{truncate}_k(U, \Sigma, V^T)$
\end{algorithmic}
\end{algorithm}